\documentclass[11pt]{article}

\usepackage[margin=1in]{geometry}
\usepackage{amsmath}
\usepackage{hyperref}
\usepackage{enumitem}
\usepackage{booktabs}
\usepackage{titlesec}
\usepackage{setspace}
\usepackage{graphicx}

\titleformat{\section}{\large\bfseries}{}{0em}{}
\setstretch{1.1}

\title{\textbf{Systemic Risk Exchange (SRX)}\\
Global Infrastructure for Portfolio-Specific Systemic Risk Hedging}
\author{}
\date{}

\begin{document}

\maketitle

\begin{abstract}
Systemic Risk Exchange (SRX) is a prototype financial infrastructure platform designed to measure systemic market stress, model portfolio-specific drawdowns, estimate crash protection premiums, and simulate clearinghouse resilience during extreme financial events.

Unlike traditional hedging instruments such as index options or volatility derivatives, SRX reduces basis risk by constructing synthetic portfolio risk indices tailored to each participant’s asset exposure.

The platform integrates systemic stress detection, cross-asset analytics, crash protection pricing, dynamic circuit breakers, and a clearinghouse-style default waterfall. This paper describes the architecture of the prototype system and outlines a staged development roadmap from research prototype to potential institutional financial infrastructure.
\end{abstract}

\section{Introduction}

Systemic financial crises remain among the most difficult risks to hedge efficiently. Investors commonly rely on broad market instruments such as index put options or volatility derivatives to protect portfolios against extreme drawdowns.

While useful, these instruments suffer from structural limitations:

\begin{itemize}
\item significant basis risk relative to portfolio exposures
\item reduced liquidity during periods of systemic stress
\item absence of integrated infrastructure for transferring systemic risk
\end{itemize}

SRX proposes a framework that measures systemic stress directly and links those measurements to portfolio-specific protection pricing and clearinghouse-style risk safeguards.

\section{System Architecture Overview}

The SRX prototype consists of five integrated system layers:

\begin{enumerate}
\item Portfolio Risk Mapping Engine (PRI)
\item Systemic Risk Modeling Engine (SRS)
\item Global Systemic Risk Benchmark (GSRI)
\item Crash Protection Pricing Engine
\item Clearinghouse Risk Management Framework
\end{enumerate}

\begin{figure}[h]
\centering
\fbox{\parbox{12cm}{
\centering
Architecture Diagram Placeholder

Portfolio Data → PRI Engine → SRS / GSRI Engine → Pricing Engine → Clearinghouse Layer
}}
\caption{High-Level SRX System Architecture}
\end{figure}

The current implementation is a research prototype developed in Python and designed to demonstrate the architecture and economic logic of a systemic risk platform.

\section{Portfolio Risk Mapping Engine}

The portfolio decomposition layer converts heterogeneous investment portfolios into standardized systemic risk representations.

Participants submit portfolios containing:

\begin{itemize}
\item equities
\item fixed income securities
\item derivatives exposures
\item commodities
\item digital assets
\item structured financial products
\end{itemize}

Each asset is mapped to systemic risk factors:

\begin{verbatim}
RiskVector = [
    EquityBeta,
    InterestRateDuration,
    CreditSpreadSensitivity,
    CommodityExposure,
    FXExposure,
    VolatilitySensitivity
]
\end{verbatim}

The system constructs a \textbf{Portfolio Risk Index (PRI)}:

\[
PRI_t = \sum_{i=1}^{N} w_i R_i(t)
\]

where

\begin{itemize}
\item $w_i$ represents portfolio weights
\item $R_i(t)$ represents returns of systemic risk factors
\end{itemize}

In the prototype implementation, PRI is represented as a normalized time series beginning at 100 and evolving through time.

\section{Systemic Risk Modeling Engine}

SRX quantifies systemic stress using multiple interacting market signals.

\subsection{Data Infrastructure and Market Oracles}

The prototype ingests historical market data from public financial data providers. In a production environment, the platform would incorporate institutional-grade market data feeds including:

\begin{itemize}
\item exchange trade and order book feeds
\item corporate credit spread data
\item volatility surface data
\item cross-asset macroeconomic indicators
\end{itemize}

These inputs act as the system's \textbf{market oracle layer}, providing the real-time signals used for systemic stress measurement.

\subsection{Training the Systemic Risk Model}

The weighting function used in the Systemic Risk Score is calibrated using historical crisis data and simulated stress environments.

Synthetic data generation techniques can be used to create simulated market environments representing extreme liquidity shocks, volatility regime transitions, and cross-asset contagion events. These simulations help estimate the sensitivity of systemic indicators and improve the robustness of the stress detection framework.

\subsection{Core Systemic Signals}

The Systemic Risk Score (SRS) is constructed from four principal signals.

\textbf{Volatility Signal ($V_t$)}

\begin{itemize}
\item rolling cross-asset volatility regimes
\end{itemize}

\textbf{Liquidity Signal ($L_t$)}

Liquidity stress is measured using the \textbf{Amihud Illiquidity Ratio}, defined as:

\[
Amihud_t = \frac{|R_t|}{DollarVolume_t}
\]

This metric captures price impact per unit of trading volume and provides a direct proxy for market liquidity deterioration. During systemic events, price movements tend to become disproportionately large relative to available trading volume, making this metric particularly useful for detecting liquidity breakdowns.

\textbf{Correlation Signal ($C_t$)}

\begin{itemize}
\item rolling cross-asset correlation spikes
\end{itemize}

\textbf{Stress Signal ($M_t$)}

\begin{itemize}
\item drawdown acceleration and downside shock behavior
\end{itemize}

These signals combine into the \textbf{Systemic Risk Score (SRS)}:

\[
SRS_t = f(V_t, L_t, C_t, M_t)
\]

where $f$ is a non-linear weighting function calibrated using both historical crisis data and synthetic stress simulations.

The SRS is normalized to a 0–100 scale for interpretability.

\section{Global Systemic Risk Index (GSRI)}

The \textbf{Global Systemic Risk Index (GSRI)} aggregates cross-asset stress signals into a single benchmark.

\[
GSRI_t = g(volatility, liquidity, correlations, macro\ shocks)
\]

Signals incorporated into GSRI include:

\begin{itemize}
\item equity volatility regimes
\item corporate credit stress
\item sovereign bond volatility
\item crypto market volatility
\item cross-asset correlation spikes
\end{itemize}

\begin{figure}[h]
\centering
\fbox{\parbox{12cm}{
\centering
Insert GSRI Dashboard Screenshot Here
}}
\caption{Example GSRI Visualization}
\end{figure}

GSRI serves as:

\begin{enumerate}
\item a benchmark for systemic financial stress
\item an input for protection pricing
\item a potential foundation for systemic risk derivatives
\end{enumerate}

\section{Protection Pricing Engine}

SRX estimates crash protection premiums using an actuarial-style framework:

\[
ProtectionPremium =
Notional \times BaseRate \times SRSMultiplier \times PortfolioRiskAdjustment \times LiquidityAdjustment
\]

Pricing inputs include:

\begin{itemize}
\item portfolio volatility
\item systemic stress level
\item liquidity conditions
\item notional exposure
\end{itemize}

Premiums increase nonlinearly as systemic stress rises.

\section{Stress Testing Layer}

SRX includes scenario-based stress testing:

\begin{itemize}
\item 2008 financial crisis
\item COVID-19 crash
\item interest rate shock
\item liquidity freeze
\item crypto contagion
\end{itemize}

For each scenario the system estimates:

\begin{itemize}
\item projected portfolio drawdown
\item stressed volatility
\item worst asset contributor
\end{itemize}

\section{Central Clearinghouse Framework}

SRX is conceptually modeled as a central counterparty clearinghouse.

\begin{verbatim}
Buyer → SRX Clearinghouse → Liquidity Provider Pools
\end{verbatim}

Losses are absorbed through a structured default waterfall:

\begin{enumerate}
\item Defaulter Initial Margin
\item Defaulter Default Fund Contribution
\item Defaulter Contingent Capital
\item Mutualized Default Fund
\item SRX Clearinghouse Capital
\end{enumerate}

\section{Dynamic Circuit Breakers}

SRX includes automated risk controls triggered by systemic stress.

\begin{center}
\begin{tabular}{|c|p{9cm}|}
\hline
\textbf{SRS Level} & \textbf{Automatic Risk Control Action} \\
\hline
Normal & Standard operation \\
\hline
Elevated & Margin requirements increase \\
\hline
Restricted & New protection contracts restricted \\
\hline
Critical & Circuit breakers activate and contingent capital reserves may be called \\
\hline
\end{tabular}
\end{center}

These controls reduce the probability of cascading defaults during extreme market stress.

\section{Institutional Risk Operating System}

The SRX prototype also functions as a risk analytics platform.

Capabilities include:

\begin{itemize}
\item systemic risk monitoring
\item portfolio stress testing
\item crash protection pricing
\item clearinghouse stress simulation
\end{itemize}

This analytics layer represents the most practical initial deployment path.

\section{Development Roadmap}

SRX development follows a staged roadmap designed to reduce execution risk.

\begin{enumerate}
\item Research prototype and methodology validation
\item Public systemic risk analytics platform
\item Institutional pilot deployments
\item Bilateral protection contracts
\item Potential evolution into regulated market infrastructure
\end{enumerate}

This staged approach allows SRX to begin as a risk analytics platform before evolving toward more complex financial infrastructure.

\section{Conclusion}

Systemic risk remains one of the least efficiently transferred risks in global financial markets. Existing hedging tools provide incomplete protection and often fail during extreme market stress.

SRX proposes a framework that combines systemic risk measurement, portfolio-specific protection pricing, and clearinghouse-style safeguards. While currently implemented as a research prototype, the architecture demonstrates how a next-generation systemic risk platform could operate and evolve toward institutional financial infrastructure.

\end{document}